% !PRINTED 629
\weberpage{629}
%<<CONT
and if we apply this to every factor of $F_n(t)$, and observe
that $p^{h+1}$ runs through the same sequence of numbers with
respect to the modulus $m$ as $p^h$ does, it follows:
\[
[F_n(t)]^p \equiv F_n(t^p) \ (\mathrm{mod}\ p),
\]
and this congruence then naturally holds also for the modulus $\mathfrak p$.

But this is the criterion for $F_n(t)$ being congruent, with
respect to the modulus $\mathfrak p$, to a polynomial $P_n(t)$ with integral
rational numerical coefficients (\S.\,150, 4.); thus
\weq{4}{F_n(t) \equiv P_n(t) \ (\mathrm{mod}\ \mathfrak p).}

If we substitute this in (3), there results in the first place
a congruence with respect to the modulus $\mathfrak p$, which however,
since it is a congruence between \gesperrt{rational} functions, must also
hold for the modulus $p$; thus
\weq{5}{f(t) \equiv P_{\xi_1}(t)\ P_{\xi_2}(t)\ \ldots\ P_{\xi_e}(t) \ (\mathrm{mod}\ p).}

By the definition (1), (4) we have
\weq{6}{P_1(r) \equiv 0 \ (\mathrm{mod}\ \mathfrak p),}
and all the roots of this congruence are
\weq{7}{r,\ r^p,\ \ldots,\ r^{p^{f-1}}.}

In the same way the congruence
\[
P_n(t) \equiv 0 \ (\mathrm{mod}\ \mathfrak p)
\]
has the roots
\[
r^n,\ r^{np},\ \ldots,\ r^{np^{f-1}},
\]
and these roots, if one lets $n$ run through the system of numbers
$\xi_1, \xi_2, \ldots, \xi_e$, are all incongruent with respect to the
modulus $\mathfrak p$.

If, in the congruence
\[
F_1(t) \equiv P_1(t) \ (\mathrm{mod}\ \mathfrak p)
\]
one makes the substitution $(r, r^n)$, then $F_1(t)$ passes into $F_n(t)$,
and $\mathfrak p$ into the conjugate ideal $\mathfrak p_n$, while $P_1(t)$, as a rational
form, remains unchanged. We therefore have
\weq{8}{F_n(t) \equiv P_1(t) \ (\mathrm{mod}\ \mathfrak p_n),}
and if we set $t=r$ in this, it follows:
\weq{9}{F_n(r) \equiv P_1(r) \ (\mathrm{mod}\ \mathfrak p_n),}
$F_n(r)$ is, however, if $n$ is not contained in $\mathfrak A$, relatively
prime to $p$, and hence not divisible by $\mathfrak p_n$.

By (9), then, $P_1(r)$ is divisible by $\mathfrak p_1$, but by none of the
%CONT>>
